I need to reconstruct and rigorously prove the unique ergodicity result for the stochastic skew heat equation, filling in the technical details. Let me organize the complete argument.

\section{Setup and Main Theorem}

\begin{definition}\label{def:setup}
Let $D:=[0,1]$ and let $p_t(x,y)$ denote the Dirichlet heat kernel on $D$, i.e.\ the kernel of the semigroup $P_t=e^{t\partial_x^2/2}$ acting on functions vanishing at the endpoints. (We normalize so that $\partial_t p=\tfrac12\partial_x^2 p$; the precise constant is immaterial.) The kernel admits the spectral representation
\[
p_t(x,y)=2\sum_{n=1}^\infty e^{-\frac12 n^2\pi^2 t}\sin(n\pi x)\sin(n\pi y),\qquad x,y\in D,\ t>0,
\]
and $P_tf(x)=\int_D p_t(x,y)f(y)\,dy$. We have the standard bounds: $p_t(x,y)\ge 0$, $\int_D p_t(x,y)\,dy\le 1$, and there is $c>0$ with
\begin{equation}\label{eq:hk-gauss}
p_t(x,y)\le c\,t^{-1/2}\exp\!\Big(-\frac{(x-y)^2}{2t}\Big),\qquad t\in(0,1].
\end{equation}
\end{definition}

Let $W$ be a space--time white noise on $\R_+\times D$ over a filtered probability space $(\Omega,\F,(\F_t)_{t\ge0},\P)$, with $(\F_t)$ the (augmented) filtration generated by $W$. Define the stochastic convolution
\[
V(t,x):=\int_0^t\int_D p_{t-r}(x,y)\,W(dr,dy),\qquad V_t(x):=V(t,x).
\]
Throughout, $h\colon\R\to\R$ is smooth with $\|h\|_{\C^{-1}}\le 1$, where for a bounded continuous $b$,
\[
\|b\|_{\C^{-1}}:=\sup_{\sigma>0}\sigma\,\|G_{\sigma^2}b\|_{L_\infty(\R)}+\|b\|_{L_\infty},
\quad
G_{\sigma^2}b(x):=\int_\R \frac{1}{\sqrt{2\pi}\,\sigma}e^{-\frac{(x-y)^2}{2\sigma^2}}b(y)\,dy.
\]
We use the standard heat-flow characterization of negative Hölder norms: $\|G_{\sigma^2}b\|_{\C^k(\R)}\le C\,\|b\|_{\C^{-1}}\sigma^{-1-k}$ for $k\ge 0$, and for $\gamma\in(0,1)$, $\|G_{\sigma^2}b\|_{\C^\gamma(\R)}\le C\|b\|_{\C^{-1}}\sigma^{-1-\gamma}$. We will use:
\begin{equation}\label{eq:gauss-bound}
\|G_{\sigma^2}b\|_{L_\infty}\le \|b\|_{\C^{-1}}\sigma^{-1},\quad
\|G_{\sigma^2}b\|_{\C^1}\le C\|b\|_{\C^{-1}}\sigma^{-2},\quad
\|G_{\sigma^2}b\|_{\C^{2/3}}\le C\|b\|_{\C^{-1}}\sigma^{-5/3}.
\end{equation}

We consider on $C_0(D):=\{f\in C(D):f(0)=f(1)=0\}$ (with sup-norm $\|\cdot\|_\infty=\|\cdot\|_\infty$, written $\|\cdot\|_\dis$) the mild solutions
\begin{align}
u_t(x)&=P_tu_0(x)+\int_0^t\!\!\int_D p_{t-r}(x,y)h(u_r(y))\,dy\,dr+V_t(x),\label{eq:u}\\
v_t(x)&=P_tv_0(x)+\int_0^t\!\!\int_D p_{t-r}(x,y)h(v_r(y))\,dy\,dr+V_t(x),\label{eq:v}\\
\wt v_t(x)&=P_tv_0(x)+\int_0^t\!\!\int_D p_{t-r}(x,y)\big(h(\wt v_r(y))+\lambda_r(u_r(y)-\wt v_r(y))\big)\,dy\,dr+V_t(x),\label{eq:vt}
\end{align}
where $\lambda\colon[0,\infty)\to\R_+$ is deterministic and piecewise constant: $\lambda_t=\lambda_n$ for $t\in[n,n+1)$. Since $h$ is smooth and bounded together with its derivative (indeed we ultimately approximate a general continuous $h$ by mollifications $h=g_{1/m}$), \eqref{eq:u}--\eqref{eq:vt} are well-posed with continuous adapted solutions taking values in $C_0(D)$; this is classical for one-dimensional SPDEs with Lipschitz nonlinearity. We write $\PP^h_n$ for the time-$n$ Markov semigroup of \eqref{eq:u}.

We will prove:

\begin{theorem}\label{thm:main}
The Markov semigroup $\PP^h$ associated with the stochastic skew heat equation \eqref{eq:u} admits at most one invariant probability measure on $C_0(D)$.
\end{theorem}

The proof occupies the rest of the document. We first establish the stochastic-sewing estimates (\Cref{L26} and its corollaries), then the a priori Hölder bound (\Cref{l:apriori}), then \Cref{L:main}, and finally the contraction/coupling argument.

\section{Notation for stochastic sewing}

For $0\le S<T$ set $\Delta_{[S,T]}:=\{(s,t):S\le s<t\le T\}$ and $\Delta^3_{[S,T]}:=\{(s,u,t):S\le s<u<t\le T\}$. For $A\colon\Delta_{[S,T]}\to L_m(\Omega)$ put $\delta A_{s,u,t}:=A_{s,t}-A_{s,u}-A_{u,t}$. Denote $\E^s[\cdot]=\E[\cdot|\F_s]$, $\|\cdot\|=\|\cdot\|_{L_m(\Omega)}$. We freely use the conditional Gaussian computation: since $V(r,y)-\E^s V(r,y)$ is, conditionally on $\F_s$, centered Gaussian with variance $\sigma^2(s,r,y):=\E\big[(V(r,y)-\E^s V(r,y))^2\big]$ and independent of $\F_s$, we have for any bounded measurable $b$ and any $\F_s$-measurable $c$,
\begin{equation}\label{eq:condgauss}
\E^s\big[b(V(r,y)+c)\big]=\big(G_{\sigma^2(s,r,y)}b\big)\big(\E^s V(r,y)+c\big).
\end{equation}

We recall L\^e's stochastic sewing lemma in the form we use.

\begin{lemma}[Stochastic Sewing, {\cite{LeSSL}}]\label{lem:SSL}
Let $m\ge2$ and $(s,t)\mapsto A_{s,t}\in L_m(\Omega)$ with $A_{s,t}$ $\F_t$-measurable. Suppose there are $\Gamma_1,\Gamma_2\ge0$, $\varepsilon_1>1/2$, $\varepsilon_2>1$ with
\[
\|A_{s,t}\|_{L_m}\le \Gamma_1(t-s)^{\varepsilon_1},\qquad
\|\E^s\delta A_{s,u,t}\|_{L_m}\le \Gamma_2(t-s)^{\varepsilon_2}.
\]
Then there is a unique (up to modification) adapted process $\A_t$ with $\A_S=0$ and
\[
\|\A_t-\A_s-A_{s,t}\|_{L_m}\le C(\Gamma_1(t-s)^{\varepsilon_1}+\Gamma_2(t-s)^{\varepsilon_2}),\quad (s,t)\in\Delta_{[S,T]},
\]
where $C=C(m,\varepsilon_1,\varepsilon_2)$. Moreover if $\A_t=\lim \sum A_{t_i,t_{i+1}}$ along partitions, then $\|\A_t-\A_s\|_{L_m}\le C(\Gamma_1(t-s)^{\varepsilon_1}+\Gamma_2(t-s)^{\varepsilon_2})$.
\end{lemma}

\section{Lower bound for the conditional variance}

\begin{lemma}\label{L:fkb}
There is $C>0$ such that for $(s,t)\in\Delta_{[0,1]}$, $x\in D$,
\[
\sigma^2(s,t,x)\ge C\big((1-x)\wedge x\wedge (t-s)^{1/2}\big).
\]
\end{lemma}

\begin{proof}
By the Markov / independent-increment structure of $V$ (the increment of the noise on $[s,t]$ is independent of $\F_s$),
\[
V(t,x)-\E^s V(t,x)=\int_s^t\int_D p_{t-r}(x,y)\,W(dr,dy)+\text{($\F_s$-measurable, mean zero given $\F_s$ correction)}.
\]
More precisely, writing $V(t,x)=\int_0^t\int_D p_{t-r}(x,y)W(dr,dy)$, the part of the Wiener integral over $r\in[s,t]$ is independent of $\F_s$, while the part over $r\in[0,s]$ is $\F_s$-measurable. Hence
\[
\sigma^2(s,t,x)\ge \E\Big[\Big(\int_s^t\int_D p_{t-r}(x,y)W(dr,dy)\Big)^2\Big]=\int_s^t\int_D p_{t-r}(x,y)^2\,dy\,dr=\int_0^{t-s}p_{2\rho}(x,x)\,\frac{d\rho}{1}.
\]
Here we used the Chapman--Kolmogorov identity $\int_D p_a(x,y)p_a(x,y)\,dy=p_{2a}(x,x)$. Thus
\[
\sigma^2(s,t,x)\ge \tfrac12\int_0^{2(t-s)}p_{\rho}(x,x)\,d\rho\ge \tfrac12\int_0^{(t-s)}p_\rho(x,x)\,d\rho.
\]
Assume WLOG $x\le 1-x$. By the spectral formula and a comparison with the whole-line/half-line heat kernel, there is $c_0>0$ such that
\[
p_\rho(x,x)\ge \frac{c_0}{\sqrt\rho}\qquad\text{whenever } \rho\le x^2 \text{ and } \rho\le \tfrac1{16}.
\]
Indeed, for the Dirichlet kernel on $[0,1]$, $p_\rho(x,x)=q_\rho(x,x)-(\text{reflection terms})$ where $q$ is the whole-line kernel; the leading reflection term is bounded by $c\,\rho^{-1/2}e^{-2x^2/\rho}$, which for $\rho\le x^2$ is $\le c\,\rho^{-1/2}e^{-2}$, while $q_\rho(x,x)=(2\pi\rho)^{-1/2}$; the further reflections are summable and exponentially small. Choosing constants gives the claim. Therefore
\[
\sigma^2(s,t,x)\ge \tfrac12\int_0^{(t-s)\wedge x^2\wedge\frac1{16}}\frac{c_0}{\sqrt\rho}\,d\rho
= c_0\sqrt{(t-s)\wedge x^2\wedge\tfrac1{16}}\ge C\big((t-s)^{1/2}\wedge x\big),
\]
which gives the stated bound (recalling $x\le 1-x$).
\end{proof}

\section{Spatial integral bounds for negative powers of $\sigma$}

\begin{lemma}\label{L:skb}
For $\rho\in(-2,0]$ and $\mu\in(0,1]$ there is $C=C(\rho,\mu)$ such that for $a\in[0,1]$, $\kappa\in(0,1]$, $x,z\in D$,
\begin{align*}
&\int_D p_\kappa(x,y)\big(y\wedge(1-y)\wedge a\big)^\rho\,dy\le C\,a^\rho\Big(1+\frac{a^2}{\kappa}\Big),\\
&\int_D|p_\kappa(x,y)-p_\kappa(z,y)|\big(y\wedge(1-y)\wedge a\big)^\rho\,dy\le C\,|x-z|^\mu\kappa^{-\mu/2}a^\rho\Big(1+\frac{a^2}{\kappa}\Big).
\end{align*}
\end{lemma}

\begin{proof}
Write $w(y):=(y\wedge(1-y)\wedge a)^\rho$. Since $\rho\le0$, $w(y)\le a^\rho$ for $y$ with $y\wedge(1-y)\ge a$, and $w(y)=(y\wedge(1-y))^\rho$ otherwise. The function $w$ is integrable since $\rho>-2>-1$: $\int_0^a y^\rho dy=\frac{a^{\rho+1}}{\rho+1}$.

\emph{First bound.} Split $D=\{y\wedge(1-y)\ge a\}\cup\{y\wedge(1-y)<a\}$. On the first set, $w(y)\le a^\rho$, so its contribution is $\le a^\rho\int_D p_\kappa(x,y)\,dy\le a^\rho$. On the set $\{y<a\}$ (similarly $\{1-y<a\}$), using $p_\kappa(x,y)\le c\kappa^{-1/2}$ from \eqref{eq:hk-gauss},
\[
\int_0^a p_\kappa(x,y)y^\rho\,dy\le c\,\kappa^{-1/2}\int_0^a y^\rho\,dy=\frac{c}{\rho+1}\kappa^{-1/2}a^{\rho+1}
=\frac{c}{\rho+1}a^\rho\cdot\frac{a}{\sqrt\kappa}\le \frac{c}{\rho+1}a^\rho\Big(1+\frac{a^2}{\kappa}\Big),
\]
using $a/\sqrt\kappa\le 1+a^2/\kappa$. Summing both contributions gives the first bound.

\emph{Second bound.} By the spectral formula, $|p_\kappa(x,y)-p_\kappa(z,y)|=\big|2\sum_n e^{-\frac12 n^2\pi^2\kappa}(\sin n\pi x-\sin n\pi z)\sin n\pi y\big|$. Using $|\sin n\pi x-\sin n\pi z|\le \min(2, n\pi|x-z|)\le C|x-z|^\mu n^\mu$ for $\mu\in(0,1]$, we get the pointwise gradient-type bound
\[
|p_\kappa(x,y)-p_\kappa(z,y)|\le C|x-z|^\mu\sum_{n\ge1}e^{-\frac12 n^2\pi^2\kappa}n^\mu\le C|x-z|^\mu\kappa^{-\frac12-\frac\mu2},
\]
where we used the elementary estimate $\sum_{n\ge1}e^{-cn^2\kappa}n^\mu\le C\kappa^{-(1+\mu)/2}$ valid for $\mu>-1$, $\kappa\in(0,1]$. Hence, exactly as above,
\[
\int_D|p_\kappa(x,y)-p_\kappa(z,y)|w(y)\,dy\le C|x-z|^\mu\kappa^{-\frac12-\frac\mu2}\int_D w(y)\,dy.
\]
Now $\int_D w\le \int_{\{y\wedge(1-y)\ge a\}}a^\rho dy+2\int_0^a y^\rho dy\le a^\rho+\frac{2}{\rho+1}a^{\rho+1}\le C a^\rho$ (as $a\le1$). Thus the integral is $\le C|x-z|^\mu\kappa^{-\frac12-\frac\mu2}a^\rho$. We must produce the factor $\kappa^{-\mu/2}a^\rho(1+a^2/\kappa)$. Write $\kappa^{-1/2}=\kappa^{-1/2}$; since $1\le 1+a^2/\kappa$ trivially gives the weaker statement, but we need to remove one power of $\kappa^{-1/2}$. Refine: split as in the first bound. On $\{y\wedge(1-y)\ge a\}$, $w\le a^\rho$ and we use the gradient bound to get contribution $\le C|x-z|^\mu\kappa^{-\frac12-\frac\mu2}a^\rho$. But we can instead use that $|p_\kappa(x,\cdot)-p_\kappa(z,\cdot)|$ integrates to $\le C|x-z|^\mu\kappa^{-\mu/2}$ over $D$:
\[
\int_D|p_\kappa(x,y)-p_\kappa(z,y)|\,dy\le C|x-z|^\mu\kappa^{-\mu/2},
\]
which follows by interpolation between $\int_D|p_\kappa(x,y)-p_\kappa(z,y)|dy\le 2$ and $\int_D|\partial_x p_\kappa(x,y)|dy\le C\kappa^{-1/2}$ (the latter from $\sum_n e^{-cn^2\kappa}n\le C\kappa^{-1}$ and the $\sin'$ structure; standard). Then on $\{y\wedge(1-y)\ge a\}$ the contribution is $\le a^\rho\cdot C|x-z|^\mu\kappa^{-\mu/2}$. On $\{y<a\}$ we use the pointwise gradient bound $|p_\kappa(x,y)-p_\kappa(z,y)|\le C|x-z|^\mu\kappa^{-\frac12-\frac\mu2}$ and $\int_0^a y^\rho dy=\frac{a^{\rho+1}}{\rho+1}$:
\[
C|x-z|^\mu\kappa^{-\frac12-\frac\mu2}a^{\rho+1}=C|x-z|^\mu\kappa^{-\mu/2}a^\rho\cdot\frac{a}{\sqrt\kappa}\le C|x-z|^\mu\kappa^{-\mu/2}a^\rho\Big(1+\frac{a^2}{\kappa}\Big).
\]
Combining gives the second bound.
\end{proof}

\begin{corollary}\label{c:key}
For $\rho\in(-4,0]$, $\mu\in(0,1]$ there is $C=C(\rho,\mu)$ such that for $x,z\in D$, $(s,t)\in\Delta_{[0,1]}$, $\kappa\in[t-s,1]$,
\[
\int_D p_\kappa(x,y)\,\sigma(s,t,y)^\rho\,dy\le C(t-s)^{\rho/4},
\quad
\int_D|p_\kappa(x,y)-p_\kappa(z,y)|\,\sigma(s,t,y)^\rho\,dy\le C|x-z|^\mu(t-s)^{\frac\rho4-\frac\mu2}.
\]
\end{corollary}

\begin{proof}
By \Cref{L:fkb}, $\sigma(s,t,y)^2\ge C(y\wedge(1-y)\wedge(t-s)^{1/2})$, hence for $\rho\le0$,
\[
\sigma(s,t,y)^\rho\le C^{\rho/2}\big(y\wedge(1-y)\wedge(t-s)^{1/2}\big)^{\rho/2}.
\]
Apply \Cref{L:skb} with exponent $\rho/2\in(-2,0]$ and $a:=(t-s)^{1/2}\le 1$. Since $\kappa\ge t-s=a^2$, we have $1+a^2/\kappa\le 2$. The first bound gives
\[
\int_D p_\kappa(x,y)\sigma(s,t,y)^\rho dy\le C\,a^{\rho/2}\cdot 2=C(t-s)^{\rho/4}.
\]
The second gives $\le C|x-z|^\mu\kappa^{-\mu/2}a^{\rho/2}\cdot 2$. Since $\kappa\ge t-s$, $\kappa^{-\mu/2}\le (t-s)^{-\mu/2}$, so the bound is $\le C|x-z|^\mu(t-s)^{-\mu/2}(t-s)^{\rho/4}=C|x-z|^\mu(t-s)^{\rho/4-\mu/2}$.
\end{proof}

\section{The key sewing lemma}

We use the seminorm
\[
[f]_{\X}=[f]_{\C^{\tau,0}L_m([S,T])}:=\sup_{(s,t)\in\Delta_{[S,T]}}\sup_{x\in D}\frac{\|f_t(x)-\E^s f_t(x)\|_{L_m}}{(t-s)^\tau},\quad \tau=\tfrac34,
\]
and $\X:=\C^{3/4,0}L_m([S,T])$, and the norm $\|f\|_{\C^{0,0}L_m([S,T])}=\sup_{s\in[S,T],x\in D}\|f_s(x)\|_{L_m}$.

\begin{lemma}\label{L26}
Let $m\ge2$, $\mu\in(0,\tfrac16)$. There is $C=C(\mu,m)>0$ such that for any bounded continuous $b\colon\R\to\R$, any continuous adapted $\phi,\psi\colon\Omega\times[0,1]\times D\to\R$, any $x,z\in D$, any $0\le S\le T\le T_0\le1$ with $T-S\le T_0-T$, setting $\Gamma:=C\|b\|_{\C^{-1}}\|w\|_{L_\infty([S,T])}$ for a continuous $w\colon[0,1]\to\R$,
\[
\A^\phi_t(x):=\int_0^t\!\!\int_D w_r\,p_{T_0-r}(x,y)\,b(V_r(y)+\phi_r(y))\,dr\,dy,\quad
\A^{\phi,\psi}_{s,t}(x):=(\A^\phi-\A^\psi)_t(x)-(\A^\phi-\A^\psi)_s(x),
\]
we have, with $\X=\C^{3/4,0}L_m([S,T])$,
\begin{align}
&\|\A^\phi_T(x)-\A^\phi_S(x)\|\le \Gamma(T-S)^{3/4}\big(1+[\phi]_\X(T-S)^{1/2}\big),\label{eq:mr1}\\
&\|\A^{\phi,\psi}_{S,T}(x)\|\le \Gamma(T-S)^{7/12}\|\psi-\phi\|_{\C^{0,0}L_m([S,T])}^{2/3}+\Gamma(T-S)^{5/4}\big([\phi]_\X+[\psi]_\X\big),\label{eq:mr2}\\
&\|\A^{\phi,\psi}_{S,T}(x)-\A^{\phi,\psi}_{S,T}(z)\|\le \Gamma|x-z|^\mu(T-S)^{\frac7{12}-\frac\mu2}\|\psi-\phi\|_{\C^{0,0}L_m}^{2/3}\nonumber\\
&\hspace{18em}+\Gamma|x-z|^\mu(T-S)^{\frac54-\frac\mu2}\big([\phi]_\X+[\psi]_\X\big).\label{eq:mr3}
\end{align}
\end{lemma}

\begin{proof}
Fix $S,T,T_0,x$ as stated. Throughout, when applying \Cref{c:key} with $\kappa=T_0-r$ for $r\in[s,t]\subseteq[S,T]$, we have $\kappa=T_0-r\ge T_0-T\ge T-S\ge t-s$, so the hypothesis $\kappa\ge t-s$ of \Cref{c:key} holds.

\smallskip
\emph{Proof of \eqref{eq:mr1}.} Define the germ
\[
A^\phi_{s,t}:=\E^s\int_s^t\!\!\int_D w_r p_{T_0-r}(x,y)\,b(V_r(y)+\E^s\phi_r(y))\,dr\,dy,\quad (s,t)\in\Delta_{[S,T]}.
\]
\emph{(i) Bound on $A^\phi_{s,t}$.} By \eqref{eq:condgauss} (with $c=\E^s\phi_r(y)$, $\F_s$-measurable),
\[
\E^s b(V_r(y)+\E^s\phi_r(y))=\big(G_{\sigma^2(s,r,y)}b\big)\big(\E^s V_r(y)+\E^s\phi_r(y)\big),
\]
hence by Fubini, \eqref{eq:gauss-bound}, and $|w_r|\le\|w\|_{L_\infty}$,
\[
|A^\phi_{s,t}|\le \|w\|_{L_\infty}\int_s^t\!\!\int_D p_{T_0-r}(x,y)\|G_{\sigma^2(s,r,y)}b\|_{L_\infty}\,dy\,dr
\le \|w\|_{L_\infty}\|b\|_{\C^{-1}}\int_s^t\!\!\int_D p_{T_0-r}(x,y)\sigma(s,r,y)^{-1}\,dy\,dr.
\]
By \Cref{c:key} with $\rho=-1\in(-4,0]$, $\kappa=T_0-r$, applied with the pair $(s,r)$ in place of $(s,t)$ (and $\sigma(s,r,y)\ge\sigma$ lower bound from \Cref{L:fkb}; note $\kappa=T_0-r\ge r-s$ as required):
\[
\int_D p_{T_0-r}(x,y)\sigma(s,r,y)^{-1}\,dy\le C(r-s)^{-1/4}.
\]
Therefore $|A^\phi_{s,t}|\le C\|w\|_{L_\infty}\|b\|_{\C^{-1}}\int_s^t(r-s)^{-1/4}\,dr\le \Gamma(t-s)^{3/4}$. This is deterministic, so $\|A^\phi_{s,t}\|_{L_m}\le\Gamma(t-s)^{3/4}$, i.e.\ $\varepsilon_1=\tfrac34>\tfrac12$.

\emph{(ii) Bound on $\E^s\delta A^\phi_{s,u,t}$.} For $(s,u,t)\in\Delta^3_{[S,T]}$,
\[
\delta A^\phi_{s,u,t}=A^\phi_{s,t}-A^\phi_{s,u}-A^\phi_{u,t}.
\]
Since $A^\phi_{s,u}$ and $A^\phi_{u,t}$ together with the part of $A^\phi_{s,t}$ over $[s,u]$ involve $\E^s$ vs $\E^u$ conditioning, the standard computation gives, taking $\E^s$ (and noting the $[s,u]$ parts cancel after applying $\E^s$ since $\E^s\E^u=\E^s$),
\[
\E^s\delta A^\phi_{s,u,t}
=\E^s\int_u^t\!\!\int_D w_r p_{T_0-r}(x,y)\Big[\E^u b(V_r(y)+\E^s\phi_r(y))-\E^u b(V_r(y)+\E^u\phi_r(y))\Big]\,dr\,dy.
\]
By \eqref{eq:condgauss},
\[
\E^u b(V_r(y)+c)=\big(G_{\sigma^2(u,r,y)}b\big)(\E^u V_r(y)+c),
\]
so with $c_1=\E^s\phi_r(y),c_2=\E^u\phi_r(y)$ (both $\F_u$-measurable), the Lipschitz bound \eqref{eq:gauss-bound} on $G_{\sigma^2(u,r,y)}b$ gives
\[
\big|\E^u b(V_r(y)+c_1)-\E^u b(V_r(y)+c_2)\big|\le \|G_{\sigma^2(u,r,y)}b\|_{\C^1}|c_1-c_2|\le C\|b\|_{\C^{-1}}\sigma(u,r,y)^{-2}|\E^s\phi_r(y)-\E^u\phi_r(y)|.
\]
Therefore
\[
|\E^s\delta A^\phi_{s,u,t}|\le \Gamma\,\E^s\int_u^t\!\!\int_D p_{T_0-r}(x,y)\sigma(u,r,y)^{-2}|\E^s\phi_r(y)-\E^u\phi_r(y)|\,dr\,dy.
\]
By Minkowski's integral inequality and the definition of $[\phi]_\X$ (with $\|\E^s\phi_r(y)-\E^u\phi_r(y)\|_{L_m}\le \|\phi_r(y)-\E^u\phi_r(y)\|_{L_m}+\|\E^s\phi_r(y)-\E^s\phi_r(y)\|=\|\E^u(\phi_r(y)-\E^s\phi_r(y))\|+\ldots$); more directly, $\|\E^s\phi_r(y)-\E^u\phi_r(y)\|_{L_m}\le \|\phi_r(y)-\E^s\phi_r(y)\|_{L_m}+\|\phi_r(y)-\E^u\phi_r(y)\|_{L_m}\le [\phi]_\X((r-s)^{3/4}+(r-u)^{3/4})\le 2[\phi]_\X(t-s)^{3/4}$. Hence
\[
\|\E^s\delta A^\phi_{s,u,t}\|_{L_m}\le \Gamma[\phi]_\X(t-s)^{3/4}\int_u^t\!\!\int_D p_{T_0-r}(x,y)\sigma(u,r,y)^{-2}\,dy\,dr.
\]
By \Cref{c:key} with $\rho=-2\in(-4,0]$ and the pair $(u,r)$ (so $\sigma(u,r,y)^{-2}$, $\kappa=T_0-r\ge r-u$),
\[
\int_D p_{T_0-r}(x,y)\sigma(u,r,y)^{-2}\,dy\le C(r-u)^{-1/2},
\]
so $\int_u^t(r-u)^{-1/2}dr=2(t-u)^{1/2}\le 2(t-s)^{1/2}$, giving
\[
\|\E^s\delta A^\phi_{s,u,t}\|_{L_m}\le \Gamma[\phi]_\X(t-s)^{3/4}(t-s)^{1/2}=\Gamma[\phi]_\X(t-s)^{5/4},
\]
i.e.\ $\varepsilon_2=\tfrac54>1$. (Here $\Gamma$ absorbed constants.)

Applying \Cref{lem:SSL} with $\Gamma_1=\Gamma$, $\Gamma_2=\Gamma[\phi]_\X$, $\varepsilon_1=\tfrac34$, $\varepsilon_2=\tfrac54$: the constructed process $\A_t$ coincides with $\A^\phi_t(x)$ because the latter is the limit of Riemann–germ sums $\sum_i A^\phi_{t_i,t_{i+1}}$. Indeed $A^\phi_{s,t}-\int_s^t\int_D w_r p_{T_0-r}b(V_r+\phi_r)$ is, after summation over a partition, $O(\text{mesh}^{1/4})$ in $L_m$ by the conditional-expectation regularity (continuity of $r\mapsto V_r,\phi_r$), so $\A^\phi_t(x)$ is the sewing of $A^\phi$. Hence
\[
\|\A^\phi_T(x)-\A^\phi_S(x)\|\le \|A^\phi_{S,T}\|+C(\Gamma_1(T-S)^{3/4}+\Gamma_2(T-S)^{5/4})\le \Gamma(T-S)^{3/4}\big(1+[\phi]_\X(T-S)^{1/2}\big),
\]
which is \eqref{eq:mr1}.

\smallskip
\emph{Proof of \eqref{eq:mr2}.} Consider the germ $A_{s,t}:=A^\phi_{s,t}-A^\psi_{s,t}$, sewn by $\A^{\phi}-\A^{\psi}$. 

\emph{(i) Bound on $A_{s,t}$.} By \eqref{eq:condgauss},
\[
A_{s,t}=\int_s^t\!\!\int_D w_r p_{T_0-r}(x,y)\Big[(G_{\sigma^2(s,r,y)}b)(\E^s V_r(y)+\E^s\phi_r(y))-(G_{\sigma^2(s,r,y)}b)(\E^s V_r(y)+\E^s\psi_r(y))\Big]dr\,dy.
\]
Using the $\C^{2/3}$ Hölder bound from \eqref{eq:gauss-bound}, $|G_{\sigma^2}b(\xi_1)-G_{\sigma^2}b(\xi_2)|\le \|G_{\sigma^2}b\|_{\C^{2/3}}|\xi_1-\xi_2|^{2/3}\le C\|b\|_{\C^{-1}}\sigma^{-5/3}|\xi_1-\xi_2|^{2/3}$,
\[
|A_{s,t}|\le \Gamma\int_s^t\!\!\int_D p_{T_0-r}(x,y)\sigma(s,r,y)^{-5/3}|\E^s\phi_r(y)-\E^s\psi_r(y)|^{2/3}\,dr\,dy.
\]
Taking $L_m$-norms (Minkowski) and using $\|\,|\E^s(\phi-\psi)_r(y)|^{2/3}\,\|_{L_m}=\|\E^s(\phi-\psi)_r(y)\|_{L_{2m/3}}^{2/3}\le \|(\phi-\psi)_r(y)\|_{L_m}^{2/3}\le \|\phi-\psi\|_{\C^{0,0}L_m([S,T])}^{2/3}$ (Jensen for conditional expectation and $L_{2m/3}\le L_m$),
\[
\|A_{s,t}\|_{L_m}\le \Gamma\|\phi-\psi\|_{\C^{0,0}L_m}^{2/3}\int_s^t\!\!\int_D p_{T_0-r}(x,y)\sigma(s,r,y)^{-5/3}\,dr\,dy.
\]
By \Cref{c:key} with $\rho=-5/3\in(-4,0]$, $\int_D p_{T_0-r}(x,y)\sigma(s,r,y)^{-5/3}dy\le C(r-s)^{-5/12}$, so $\int_s^t(r-s)^{-5/12}dr=\frac{12}{7}(t-s)^{7/12}$, giving
\[
\|A_{s,t}\|_{L_m}\le \Gamma\|\phi-\psi\|_{\C^{0,0}L_m}^{2/3}(t-s)^{7/12},
\]
i.e.\ $\varepsilon_1=\tfrac7{12}>\tfrac12$.

\emph{(ii) Bound on $\E^s\delta A_{s,u,t}$.} We have $\delta A_{s,u,t}=\delta A^\phi_{s,u,t}-\delta A^\psi_{s,u,t}$, hence by part (ii) of \eqref{eq:mr1}'s proof applied to each,
\[
\|\E^s\delta A_{s,u,t}\|_{L_m}\le \Gamma([\phi]_\X+[\psi]_\X)(t-s)^{5/4},
\]
i.e.\ $\varepsilon_2=\tfrac54>1$. Applying \Cref{lem:SSL} as before yields \eqref{eq:mr2}.

\smallskip
\emph{Proof of \eqref{eq:mr3}.} Consider the germ with the difference kernel $q_r(y):=p_{T_0-r}(x,y)-p_{T_0-r}(z,y)$:
\[
A^{(x,z)}_{s,t}:=\E^s\int_s^t\!\!\int_D w_r\,q_r(y)\,\big[b(V_r(y)+\E^s\phi_r(y))-b(V_r(y)+\E^s\psi_r(y))\big]\,dr\,dy,
\]
sewn by $\A^{\phi,\psi}_{s,t}(x)-\A^{\phi,\psi}_{s,t}(z)$. Repeating (i) of \eqref{eq:mr2} but with $|q_r(y)|$ in place of $p_{T_0-r}(x,y)$ and applying the second inequality of \Cref{c:key} (with $\rho=-5/3$, $\mu\in(0,\tfrac16)$, $\kappa=T_0-r\ge r-s$),
\[
\int_D|q_r(y)|\sigma(s,r,y)^{-5/3}\,dy\le C|x-z|^\mu(r-s)^{-5/12-\mu/2},
\]
so $\int_s^t(r-s)^{-5/12-\mu/2}dr\le C(t-s)^{7/12-\mu/2}$ (as $\tfrac5{12}+\tfrac\mu2<1$ since $\mu<\tfrac16$), giving
\[
\|A^{(x,z)}_{s,t}\|_{L_m}\le \Gamma|x-z|^\mu\|\phi-\psi\|_{\C^{0,0}L_m}^{2/3}(t-s)^{7/12-\mu/2},
\]
with $\varepsilon_1=\tfrac7{12}-\tfrac\mu2>\tfrac12$ (since $\mu<\tfrac16$). Similarly, repeating (ii) with the difference kernel and the second inequality of \Cref{c:key} with $\rho=-2$,
\[
\int_D|q_r(y)|\sigma(u,r,y)^{-2}\,dy\le C|x-z|^\mu(r-u)^{-1/2-\mu/2},
\]
gives $\|\E^s\delta A^{(x,z)}_{s,u,t}\|_{L_m}\le \Gamma|x-z|^\mu([\phi]_\X+[\psi]_\X)(t-s)^{5/4-\mu/2}$, with $\varepsilon_2=\tfrac54-\tfrac\mu2>1$. Applying \Cref{lem:SSL} yields \eqref{eq:mr3}.
\end{proof}

\begin{corollary}\label{c:39}
The bounds \eqref{eq:mr1}, \eqref{eq:mr2}, \eqref{eq:mr3} of \Cref{L26} hold for all $0\le S\le T\le1$ with $T_0=T$ (i.e.\ removing the restriction $T-S\le T_0-T$), with a constant $C=C(\mu,m)$.
\end{corollary}

\begin{proof}
This is the standard "taming singularities" argument \cite[Lemma 2.3]{BFG}. The only place the condition $T-S\le T_0-T$ was used is to guarantee $\kappa=T_0-r\ge t-s$ in \Cref{c:key}. For $T_0=T$ this fails near $r=T$, producing an integrable time singularity. Concretely, one applies the estimates of \Cref{L26} on dyadic subintervals $[T-2^{-k}(T-S),T-2^{-k-1}(T-S)]$ where the condition holds with $T_0=T$ (since on such an interval $[s_k,t_k]$ one has $T_0-t_k=2^{-k-1}(T-S)\ge t_k-s_k=2^{-k-1}(T-S)$), and sums the geometric series. The exponents $\tfrac34,\tfrac7{12}-\tfrac\mu2,\tfrac7{12}$ are all $>\tfrac12$ and the "second" exponents are all $>1$, so the dyadic sums converge, producing the same bounds (with a larger constant) for $T_0=T$. We omit the routine details, identical to \cite[Lemma 2.3]{BFG}.
\end{proof}

\begin{corollary}\label{c:310}
Let $m\ge2$, $\mu\in(0,\tfrac16)$. There is $C=C(\mu,m)$ such that for any bounded continuous $b$, continuous adapted $\phi,\psi$, $x,z\in D$, $0\le S\le T\le1$,
\begin{align}
&\Big\|\int_S^T\!\!\int_D w_r p_{T-r}(x,y)\big(b(V_r(y)+\phi_r(y))-b(V_r(y)+\psi_r(y))\big)dr\,dy\Big\|_{L_m}\nonumber\\
&\quad\le C\|b\|_{\C^{-1}}\|w\|_{L_\infty}(T-S)^{1/4}\|\psi-\phi\|_{\C^{0,0}L_m}+C\|b\|_{\C^{-1}}\|w\|_{L_\infty}(T-S)^{5/4}\big(1+[\phi]_\X+[\psi]_\X\big),\label{eq:c310a}\\
&\Big\|\int_S^T\!\!\int_D w_r(p_{T-r}(x,y)-p_{T-r}(z,y))\big(b(V_r(y)+\phi_r(y))-b(V_r(y)+\psi_r(y))\big)dr\,dy\Big\|_{L_m}\nonumber\\
&\quad\le C\|b\|_{\C^{-1}}\|w\|_{L_\infty}|x-z|^\mu\|\psi-\phi\|_{\C^{0,0}L_m}+C\|b\|_{\C^{-1}}\|w\|_{L_\infty}|x-z|^\mu(T-S)^{5/4-\mu/2}\big(1+[\phi]_\X+[\psi]_\X\big).\label{eq:c310b}
\end{align}
\end{corollary}

\begin{proof}
The left side of \eqref{eq:c310a} equals $\|\A^{\phi,\psi}_{S,T}(x)\|$ with $T_0=T$, so by \Cref{c:39}, \eqref{eq:mr2},
\[
\le \Gamma(T-S)^{7/12}\|\psi-\phi\|_{\C^{0,0}L_m}^{2/3}+\Gamma(T-S)^{5/4}([\phi]_\X+[\psi]_\X).
\]
Apply Young's inequality $a_1a_2\le a_1^{3/2}+a_2^3$ to $a_1=(T-S)^{1/6}\|\psi-\phi\|_{\C^{0,0}L_m}^{2/3}$, $a_2=(T-S)^{5/12}$:
\[
(T-S)^{7/12}\|\psi-\phi\|_{\C^{0,0}L_m}^{2/3}=a_1a_2\le a_1^{3/2}+a_2^3=(T-S)^{1/4}\|\psi-\phi\|_{\C^{0,0}L_m}+(T-S)^{5/4}.
\]
This yields \eqref{eq:c310a}, after absorbing the extra $(T-S)^{5/4}$ into the $(1+\ldots)$ term. For \eqref{eq:c310b}, the left side is $\|\A^{\phi,\psi}_{S,T}(x)-\A^{\phi,\psi}_{S,T}(z)\|$ with $T_0=T$; \Cref{c:39}, \eqref{eq:mr3} gives
\[
\le \Gamma|x-z|^\mu(T-S)^{\frac7{12}-\frac\mu2}\|\psi-\phi\|_{\C^{0,0}L_m}^{2/3}+\Gamma|x-z|^\mu(T-S)^{\frac54-\frac\mu2}([\phi]_\X+[\psi]_\X).
\]
Apply Young $a_1a_2\le a_1^{3/2}+a_2^3$ with $a_1=\|\psi-\phi\|_{\C^{0,0}L_m}^{2/3}$, $a_2=(T-S)^{7/12-\mu/2}$:
\[
\|\psi-\phi\|^{2/3}(T-S)^{\frac7{12}-\frac\mu2}\le \|\psi-\phi\|+(T-S)^{\frac74-\frac{3\mu}2}\le \|\psi-\phi\|+(T-S)^{5/4-\mu/2},
\]
(the last since $\tfrac74-\tfrac{3\mu}2\ge\tfrac54-\tfrac\mu2$ for $\